\documentclass[10pt,a4paper]{amsart}
\usepackage[margin=19mm]{geometry}
\usepackage{amsmath,amssymb,mathtools}
\usepackage[scr=rsfs,cal=euler]{mathalfa}
\usepackage{xfrac}
\usepackage[unicode,hidelinks,bookmarks=false]{hyperref}
\newcommand{\ie}{i.e.\ }
\newcommand{\cf}{cf.\ }
\newcommand{\resp}{respectively\ }
\newcommand{\Subsection}{\subsection}
\providecommand{\nobreakdash}{\nobreak\hbox{-}}
\setlength{\emergencystretch}{2em}
\multlinegap0pt
\usepackage{bm}
\usepackage{bbm}
\usepackage{dsfont}
\usepackage{xspace}
\usepackage{cancel}
\usepackage{mathtools}
\usepackage[shortlabels]{enumitem}
\usepackage{amsthm}
\makeatletter
\@ifundefined{theorem}{\newtheorem{theorem}{Theorem}}{}
\@ifundefined{lemma}{\newtheorem{lemma}[theorem]{Lemma}}{}
\makeatother
\title[A stability theorem for Berge Hamiltonian cycles under a min]{A stability theorem for Berge Hamiltonian cycles under a minimum degree condition}
\author{Yichen Wang, D\'aniel Gerbner, Xiamiao Zhao}
\date{Source: arXiv:2608.04603v1 (CC BY 4.0)}
\begin{document}
\maketitle

\noindent\emph{Source and licence.} A stability theorem for Berge Hamiltonian cycles under a minimum degree condition. Version arXiv:2608.04603v1 (CC BY 4.0), distributed under the Creative Commons Attribution 4.0 International licence, as recorded in the arXiv licence metadata for this exact version and shown on the abstract page https://arxiv.org/abs/2608.04603v1. Full rights notice: licenses/arxiv-2608.04603-CC-BY-4.0.txt.\par\smallskip

\noindent\textit{Editorial scope.} Counted unit: the Lemma whose source label is \texttt{lem:shadow} in the article (source0.tex lines 656--669, source label \texttt{lem:shadow}), delivered together with the article's own complete original proof of it (source0.tex lines 671--849, one \texttt{proof} environment of 179 lines). No mathematics is written, repaired or re-derived: statement and proof are the article's own text line for line. Required context retained uncounted: lem:heavy-degree (lines 594--597). Every definition, display and label of those lines is retained unchanged. Omitted from this package and not redistributed: 1--593, 598--655, 670--670, 850--989. Nothing inside a retained excerpt is omitted or altered: every retained body is delivered exactly as the article's lines are written, so no omission receipt of any kind accompanies this package. Citations: the retained excerpts carry no citation command at all, so no bibliography entry of the article is transcribed and no reference is redistributed. Reference adaptation: none was needed and none was made; every reference inside the delivered unit resolves inside it, so no unresolved reference or citation is produced. Printed slips kept literal: no printed word, symbol, hypothesis or number of the counted statement or of its proof was altered. Layout and class: the article's own class is \texttt{article} with packages including fontenc, inputenc, amsmath, amssymb, amsfonts, amsthm, authblk, natbib, enumitem, geometry, hyperref, color; the wrapper is the lane's pinned \texttt{amsart} editorial wrapper at 10pt on A4, which restates the theorem-like environments the retained text needs and adds the article's own macro definitions that the retained text needs (none), with extra wrapper packages (bm, bbm, dsfont, xspace, cancel, mathtools, enumitem). The delivered numbering is the unit's own. Assets: no retained span contains a figure environment, a graphics-inclusion command, a PDF-page inclusion, a table or a TikZ picture; no image file is copied into this package, the package declares no asset dependency and no image file is redistributed with it. Licence: version arXiv:2608.04603v1 of this article is distributed under the Creative Commons Attribution 4.0 International licence, as recorded in the arXiv licence metadata for this exact version and shown on the abstract page https://arxiv.org/abs/2608.04603v1. A full rights notice is delivered at licenses/arxiv-2608.04603-CC-BY-4.0.txt. Characters: every retained excerpt is pure ASCII, so the delivered engine, XeLaTeX with its default Latin Modern text font, reports no missing character.
\begin{lemma}\label{lem:heavy-degree}
Every vertex of $T$ has degree at least $t$ in $G_T$ and degree at least
$t-1$ in $G_h$.
\end{lemma}

\begin{lemma}[Shadow covering lemma]\label{lem:shadow}
Let $G_T$ and $G_h$ be as above.  Suppose every vertex of $T$ has degree
at least $t$ in $G_T$ and degree at least $t-1$ in $G_h$.  Then one of the
following holds.
\begin{enumerate}[label=(\roman*)]
\item $G_T$ contains a linear forest $F$ covering all vertices of $T$,
whose components have endpoints in $A$ and have no edges inside $A$.
Moreover, to the edges of $F$ not in $G_h$ 
%may be assigned in advancedistinct hyperedges  
we can assign distinct hyperedges containing them.
\item $|N_{G_T}(T)\cap A|\le 1$.
\item $G_T[T]$ is empty and $|N_{G_T}(T)\cap A|\le t$.
\end{enumerate}
\end{lemma}

\begin{proof}
    
    First we pick a linear forest $F_0$ in $G_T$ such that it consists of heavy edges inside $T$ with at most one exception as an end edge of a path component.
    Let $P_1,\dots, P_s$ be the paths in $F_0$ and assume the size of the paths is $p_1 \ge p_2 \ge \cdots \ge p_s$, such that $\sum_{i=1}^{s}p_i = t$.
    We choose $F_0$ such that $(p_1,\ldots,p_s)$ is maximized in lexicographic order.
    Subject to this, we pick $F_0$ without a light edge, if possible. If not, then there is some $i$ such that $P_i$ contains the light edge. We pick $F_0$ such that $i$ is maximized.
    %if there is a light edge in $P_i$, we pick the $F_0$ such that $i$ is maximized~(when there is no light edge, we make a deal that $i = \infty$).
By the way how we picked $F_0$, if there is a light edge in $P_i$, then $p_j < p_i$ for all $j > i$.

    \noindent
    \textbf{Case 1:} $s = t$ and $p_i = 1, i=1,2,\ldots,t$.
    
    In this case, $G_T[T]$ has no edge.
    Let $F'$ be a linear forest with endpoints in $A$, covering the maximum number of vertices in $T$, and for each covered $v_i$, at most one neighbor is light.
    If $F'$ covers all the vertices in $T$, then we are done and (i) holds, because every light pair is contained in a distinct hyperedge since $G_T[T]$ has no edge.
    
    Suppose $v_{t-1}, v_{t}$ are not covered by $F'$.
    Note that there are at least $t-1$ heavy neighbors of $v_{t-1}$ in $A$ and at least $t$ neighbors of $v_{t-1}$ in $A$.
    There are at most $t-3$ internal vertices in $A$ used in $F'$. Then we can pick a heavy neighbor of $v_{t-1}$ in $T$ avoiding the internal vertices.
    If a vertex in $F'$ is chosen, then we can pick another neighbor of $v_{t-1}$ in $A$ avoiding the internal vertices and the other endpoint of the path. This way we obtain a linear forest with endpoints in $A$, covering more vertices in $T$ than $F'$, a contradiction to the choice of $F'$.
    Otherwise, we can pick another arbitrary neighbor of $v_{t-1}$ in $A$ avoiding the internal vertices, again obtaining a linear forest with endpoints in $A$, covering more vertices in $T$ than $F'$, a contradiction to the choice of $F'$.

    Now suppose only $v_t$ is not covered.
    If $F'$ is not a single path, then there are at most $t-3$ internal vertices in $A$.
    Then we can complete the proof by a similar argument as before.
    When $F'$ contains exactly one path, there are $t-2$ internal vertices in $A$.
    The above argument only fails when the neighbor set of $v_t$ is exactly $F'\cap A$.
    Then we have a cycle in $G_T$ covering $T$ such that every $v_i$ is incident to at most one light edge.
    Assume the cycle is $v_1a_1v_2a_2\cdots v_ta_tv_1$.
    Let $A' = F' \cap A = \{a_1,a_2,\ldots,a_t\}$, then we claim that $N_{G_T}(T) = A'$ and thus we arrive to (iii). Indeed, if $v_i$ has a neighbor outside $A'$, say $a'$, without loss of generality, assume $v_ia_i$ is heavy, then we have a path $a'v_ia_i\cdots v_1a_1 \cdots v_{i-1}a_{i-1}$, a contradiction to the choice of $F'$.

     \noindent
    \textbf{Case 2:} $s=1, p_1 = t$.
    Assume the path is $v_1v_2\cdots v_t$.
    
    \textbf{Subcase 2.1:} all the edges are heavy in $F_0$. %Dani: in $F_0$??? or inside $T$???...
    
    If there is no heavy Hamiltonian cycle in $T$, then $v_1$ has a heavy neighbor $a_1 \in A$, and $v_t$ has a heavy neighbor $a_2 \neq a_1 \in A$ which leads to (i).
    Without loss of generality, assume $v_1v_2\ldots v_tv_1$ is a heavy cycle.
    
    If $v_1$ has a heavy neighbor $a_1 \in A$, 
    then we claim that for every $v_i$ with $i\neq 1$, $N(v_i)=T\cup \{a_1\}$.
    Otherwise, suppose $v_i$~($i < t$) has a neighbor $a_i \neq a_1$, then consider the heavy neighbors of $v_2$.
    If $v_2$ has a heavy neighbor $a_2 \neq a_1 \in A$, then (i) holds.
    If $v_{i+1}$ is a heavy neighbor of $v_2$, then we are done as well, because there is a heavy path with $v_1$ and $v_i$ as end vertices, and they have distinct heavy neighbors in $A$.
    Thus, according to 
    Lemma \ref{lem:heavy-degree} for vertex $v_2$, $a_1$ must be a heavy neighbor of $v_2$.
    We can continue this process and take $v_j$ as the beginning vertex for $j=2,\dots,i-1$ one by one, and we can get that $a_1$ is a heavy neighbor of $v_{i-1}$. This way we also arrive to (i).

    If $v_1$ has a heavy neighbor $a_1 \in A$, then by the previous analysis, $N(v_i) = T \cup \{a_1\}$ for all $i \neq 1$.
    If $v_1$ has no neighbors in $A\setminus \{a_1\}$, then we are done as (ii) holds.
    Otherwise, $v_1$ has a neighbor other than $a_1$ in $A$, then the heavy neighbors of each $v_i$ must be all the vertices in $T$.
    Then $G_h[T]$ is a clique.
    Note that if $v_1$ has no heavy neighbor in $A$, by the symmetry, $G_h[T]$ is a clique.
    So this is the only remaining case.

    When $G_h[T]$ is a clique, it is easy to prove that $N_{G_T}(T) \cap A$ has at most one vertex, thus (ii) holds.

    \textbf{Subcase 2.2:} there is a light edge in $F_1$, say $v_{t-1}v_{t}$. 

    We may assume $v_1v_t$ is not a heavy edge, otherwise $v_tv_1v_2\dots v_{t-1}$ is a path with both edges heavy, and we are done by the previous case.
    Since $v_1v_t$ is not heavy, $v_1$ has a heavy neighbor $a_1 \in A$.
    Note that $v_{t-1}v_{t}$ is not a heavy edge, then $v_{t}$ has at least two heavy neighbors in $A$.
    Then we are done as (i) holds.

    \noindent
    \textbf{Case 3:} $p_2 = 1$.
    
    Assume the paths are $P_1=v_1 \ldots v_{p_1}$ and the singletons $u_1, u_2,\ldots, u_q$ with $p_1 + q = t$ and $q\geq 1$.

    We shall repeatedly use the following simple observation. Let $F$ be a
linear forest whose components have endpoints in $A$, and let
$$
I(F)=\{a\in A\cap V(F): d_F(a)=2\}
$$
be the set of internal $A$-vertices of $F$. Suppose that
$u\in T\setminus V(F)$ has at least $|I(F)|+3$ neighbors in $A$.
Then $u$ can be inserted into $F$ using two of these neighbors without
creating a cycle. Indeed, after avoiding $I(F)$, at least three neighbors
of $u$ remain in $A$. Among any three such vertices, two are not the two
endpoints of the same component of $F$. Choosing such two vertices $a,b$
and adding the edges $au$ and $ub$ keeps a linear forest whose components
still have endpoints in $A$. The same statement holds with ``neighbors''
replaced by ``heavy neighbors''.

    \textbf{Subcase 3.1:} there is no light edge in $P_1$.
    
    Then each $u_i$ has at least $q+1$ heavy neighbors in $A$ and at least $q+2$ neighbors in $A$.
    $v_1, v_{p_1}$ has at least $q+1\geq 2$ heavy neighbors in $A$ and at least $q+2$ neighbors in $A$.

    First pick a heavy neighbor $a_1$ of $v_1$ in $A$, and a neighbor $a_2$ of $v_{p_1}$ in $A$ with $a_2 \neq a_1$.
    Let $F'$ be the linear forest with endpoints in $A$, covering the maximum number of vertices in $T$, containing $a_1P_1a_2$ and for each $u_i$, and there is at most one light edge in $F'$ incident to $u_i$ for some $i=1,\dots,q$.
    
    We define $I(F')$ to denote the vertices in $A\cap V(F')$ with degree two in $F'$, which is the set of internal vertices in $F'$.

    
    If $F'$ covers all the vertices in $T$, then (i) holds.
    If a vertex, say $u_q$ is not covered, then notice that $|I(F')|\leq q-2$, then we can pick a heavy neighbor $a_q$ of $u_q$ in $A$ avoiding $I(F')$.
    Since $q+2-(q-1) > 3$, we can pick another heavy neighbor of $u_q$ in $A$ avoiding $I(F')$. In the case these two neighbors are endpoints of the same path in $F'$, we can pick a third heavy neighbor. 
    %and when $u_q$ is in $F'$ avoiding the other endpoint of the path in $F'$ containing $u_q$ (which avoids the case when this extends to a cycle), a contradiction to the choice of $F'$.


    \textbf{Subcase 3.2:} $v_1v_2$ is a light edge.
    
    Then $v_1$ has at least one heavy neighbor $a_1$ in $A$.
    Note that $v_{p_1}v_1$ cannot be a heavy edge (otherwise $v_2v_3 \dots v_{p_1}v_1$ is a heavy path, which is done by the previous case), then $v_{p_1}$ has at least $q+1$ heavy neighbors in $A$.
    Pick a heavy neighbor $a_2$ of $v_{p_1}$ in $A$ with $a_2 \neq a_1$.

    For each $u_i$, when $p_1 = 2$, it cannot have $v_1$ and $v_2$ as heavy neighbors, because it will cause a heavy $P_2$, which violates the choice of $F_0$.
    When $p_1 > 2$, it cannot have $v_2$ or $v_{p_1}$ as heavy neighbors, or we can find a heavy path with length $p_1$, which also violates the choice of $F_0$.
    Then $u_i$ has at least $q+1$ heavy neighbors in $A$.

Now let us extend $u_1,\ldots,u_{q-1}$ one by one. Suppose that, before
$u_i$ is inserted, the current forest is $F_i$. Since
$|I(F_i)|\le i-1$ and $u_i$ has at least $q+1$ heavy neighbors in
$A$, for $i<q$ we have
$$
q+1 \ge |I(F_i)|+3 .
$$
By the observation above, we can choose two heavy neighbors of $u_i$ in
$A\setminus I(F_i)$ which are not the two endpoints of the same component
of $F_i$. Adding the two corresponding edges inserts $u_i$ and keeps a
linear forest with endpoints in $A$.
    
    

    When we arrive at $u_q$, let $F'$ be the current linear forest.
    If $p_1 > 2$, it fails only when $u_q$ has $v_1, v_3,v_4,\ldots, v_{p_1-1}$ as heavy neighbors.
    Since all $u_i$'s are equivalent, all the $u_i$'s have $v_1, v_3,\ldots, v_{p_1-1}$ as heavy neighbors.
    When $q > 2$ and $p_1 = 3$, we have a heavy path $u_1v_1u_2$, a contradiction.
    When $q \ge 1$ and $p_1 > 3$, we have a longer heavy path $v_2v_3\ldots v_{p_1-1}u_qv_1$, a contradiction.
    Note that it is impossible that $q = 1$ and $p_1 = 3$ since $t$ is large.

    If $p_1 = 2$, then $G_h[T]$ is empty and every vertex in $T$ has at least $t-1$ heavy neighbors in $A$.
    If there is a path of length two or a matching of size two in $G_T[T]$ such that the edges are from distinct hyperedges, then we can extend greedily to the case (i).
    So the hyperedge containing $v_1,v_2$ is the only hyperedge that contains at least two vertices in $T$.
    Without loss of generality, assume $u_q$ is not in the hyperedge since $t > r$.
    Then there are at least $t$ neighbors of $u_q$ in $A$.
    Now we can first pick a heavy neighbor $a_q$ of $u_q$ in $A$ avoiding the vertices in $I(F')$, and then pick a neighbor of $u_q$ in $A$ avoiding the internal vertices and the other endpoint of the path if $a_q$ is in $F'$, which leads to (i).


    
    \textbf{Case 4:} $s \ge 2$ and $2 \le p_1 \le t-2$.
    In this case, we would like to extend $F_0$ to a linear forest $F$ satisfying (i).
    To achieve this, we aim to extend every $P_i$ one by one using heavy neighbors.

    Start with $P_1$ with size $p_1$.
    If there is no light edge in $P_1$, then for each end point of $P_1$, the vertices in $T\setminus P_1$ cannot be heavy neighbors, so there are at least $t - p_1$ heavy neighbors in $A$.
    Since $t - p_1 \ge 2$, we can extend $P_1$ to a path $P_1'$ with both end points in $A$ using heavy neighbors.
    If there is a light edge in $P_1$, assume that $P_1 = v_1v_2\ldots v_{p_1}$ and $v_{1}v_{2}$ are light.
    Similarly, $v_1$ has at least one heavy neighbor in $A$ and $v_{p_1}$ has at least $t - p_1 \ge 2$ heavy neighbors in $A$, thus we can extend $P_1$ to a path $P_1'$ with both end points in $A$ using heavy neighbors.

    Assume that after $i$ paths in $F_0$, we arrive at a path $P_{i+1}$ with $p_{i+1} \ge 2$. 
    Let $u_1,u_2$ be the endpoints of $P_{i+1}$.
    If we haven't met a light edge before this moment, then $u_1$ and $u_2$ cannot have the previous endpoints of the paths in $F_0$ as their heavy neighbors.
    If there is a light edge in $P_{j}$ for some $j \le i$, then the internal vertex in the light edge cannot be a heavy neighbor of $u_1$ or $u_2$.
    Moreover, let $P_{i+1} = u_1u_2\ldots u_{p_{i+1}}$, then $u_{p_{i+1}}$ cannot have vertices in $P_{i+2}, \ldots, P_{s}$ as its heavy neighbors unless $u_{p_{i}}u_{p_{i+1}}$ is a light edge and $p_{i+1} \ge 3$.
    As a conclusion, one of the end points has at least $2i$ heavy neighbors in $A$ and the other has at least $2i + p_{i+2} + \ldots + p_{s}$ heavy neighbors in $A$.
    Let $F_i=P_1'\cup\dots \cup P'_{i}$, then $|I(F_i)\cap A|\leq i-1$.
    When $2i > i-1$ and $2i + p_{i+2} + \ldots + p_{s} > i-1 +2$, then we can extend $P_{i+1}$ to a path $P_{i+1}'$ with both end points in $A$ using heavy neighbors such that $P_1', \ldots, P_{i+1}'$ is a linear forest.
    It only fails when $i=1$ and $s=2$.
    In this case, $p_1 \ge 3$, otherwise we have $t \le 4$, a contradiction.
    Assume $P_1 = v_1v_2\ldots v_{p_1}$, then $v_2$ cannot be a heavy neighbor of $u_1$ or $u_2$ as well, so $u_1$ and $u_2$ both have at least three heavy neighbors in $A$ and we are done since (i) holds.
    
    Assume that after $p$ paths of length at least one and $q$ singletons in $F_0$, we arrive to a singleton $u$, and the current linear forest is $F'$. 
    Then we claim that $u$ has at least $2p+q$ heavy neighbors in $A$.
    If we haven't met a light edge before this moment, then the endpoints of these paths and the singletons, at least $2p+q$ vertices of $T$ are not heavy neighbors of $u$, thus $u$ has at least $2p+q$ heavy neighbors in $A$. 
    If there is a light edge in $P_{j}$ for some $j \le p$, let $P_j = u_1u_2\ldots u_{p_j}$ and $u_1u_2$ is light, then when $p_j = 2$, $u_1$ and $u_2$ cannot be a heavy neighbor of $u$, when $p_j > 2$, $u_2$ and $u_{p_j}$ cannot be a heavy neighbor of $u$.
    There are at most $p+q-1$ vertices in $I(F')$.
    Since $|I(F')|\le p+q-1$ and $p\ge 2$, the vertex $u$ has at least
$$
(2p+q)-(p+q-1)=p+1\ge 3
$$
heavy neighbors in $A\setminus I(F')$. By the observation above, two of
these neighbors can be chosen so that they are not the two endpoints of the
same component of $F'$. Adding the corresponding two heavy edges therefore
inserts $u$ without creating a cycle by the observation in Case 3, and the resulting graph is again a
linear forest whose components have endpoints in $A$.  
\end{proof}

\end{document}
